\documentclass[12pt,a4paper]{article}
\usepackage[utf8]{inputenc}
\usepackage[T1]{fontenc}
\usepackage[french]{babel}
\usepackage{geometry}
\usepackage{array}
\usepackage{booktabs}
\usepackage{graphicx}
\usepackage{enumitem}
\geometry{margin=1in}

\begin{document}

\begin{center}
\textbf{\Large RAPPORT D’INVESTIGATION NUMÉRIQUE}\\[1em]
\textbf{TRAVAIL N-4}\\[0.5em]
\textbf{MATRICULE : 22P056}\\
\textbf{NOMS \& PRÉNOMS : NNA FRANCIS EMMANUEL ROMEO}\\
\textbf{FILIÈRE : CIN-4}\\[2em]
\textit{Sous la supervision de : MINKA THIERRY}\\[2em]
\textbf{UNIVERSITÉ DE YAOUNDÉ I}\\
\textbf{ÉCOLE NATIONALE SUPÉRIEURE POLYTECHNIQUE DE YAOUNDÉ}\\
\textbf{DÉPARTEMENT DE GÉNIE INFORMATIQUE}\\[2em]
\textbf{RÉPUBLIQUE DU CAMEROUN}\\
\textit{PAIX – TRAVAIL – PATRIE}
\end{center}

\newpage
\tableofcontents
\newpage

\section{Introduction}
L’ordonnance de renvoi du 29 février 2024 conclut une information judiciaire menée après l’assassinat du journaliste \textbf{Martinez Zogo}. 

Le travail de l’expert judiciaire en investigation numérique a été essentiel pour relier les faits matériels aux responsabilités hiérarchiques, grâce à l’analyse des données électroniques, téléphoniques et de géolocalisation. 

Ce rapport identifie les éléments numériques effectivement utilisés et ceux qui auraient dû être produits pour appuyer les décisions du magistrat instructeur.

\section{Éléments d’investigation numérique relevés dans l’ordonnance}
\renewcommand{\arraystretch}{1.2}
\begin{tabular}{|p{4cm}|p{6cm}|p{5cm}|}
\hline
\textbf{Type d’élément numérique} & \textbf{Description} & \textbf{Utilité dans l’enquête} \\ \hline
Données de localisation téléphonique (géolocalisation) & Exploitation des données de position des téléphones de plusieurs inculpés (TONGUE NANA, LAMFU JOHNSON, DAOUDA, etc.) & Prouve la présence des suspects à Ebogo et à Soa lors des faits. Établit la chronologie des deux opérations. \\ \hline
Historique des appels téléphoniques & Examen des appels entre la victime et les inculpés (ex. appel entre Zogo et SAVOM Martin avant l’attaque). & Établit les contacts directs et la coordination avant le crime. \\ \hline
Fiche de géolocalisation (SAIWANG YVES et HEUDJI GUY SERGE) & Données extraites des systèmes de la DGRE, transmises à DANWE Justin. & Élément-clé : surveillance et repérage de la victime. \\ \hline
Conversations téléphoniques et messages WhatsApp & Échanges entre DANWE et ses subordonnés, et entre BIDJANG OB’A BIKORO et d’autres membres. & Prouve la planification, la préméditation et la coordination des missions. \\ \hline
Images de vidéosurveillance & Enregistrements des caméras autour d’Amplitude FM et du domicile de la victime. & Montre la dernière rencontre entre Zogo et BIDZONGO alias « Arthur ESSOMBA ». \\ \hline
Exploitation des téléphones placés sous scellés & Messages retrouvés sur le téléphone d’ESSOMBA (mention des « bombes » = documents sensibles). & Permet d’établir les motivations et les liens entre suspects. \\ \hline
Examen de la base de données de la DGRE & Extraction des fiches techniques internes (missions, ordres, messages). & Détermine si les opérations étaient autorisées par la hiérarchie. \\ \hline
Données bancaires électroniques & Relevés d’opérations bancaires et retraits (SAVOM Martin). & Prouve le financement matériel de l’opération et la rémunération du commando. \\ \hline
Captures d’écran de téléphones et preuves numériques imprimées & Extraits de l’historique d’appels et de messages, versés au dossier. & Servent de pièces à conviction directes. \\ \hline
\end{tabular}

\section{Éléments qui auraient dû être produits par l’expert judiciaire}
\begin{tabular}{|p{6cm}|p{8cm}|}
\hline
\textbf{Élément attendu} & \textbf{Rôle attendu dans l’analyse judiciaire} \\ \hline
Rapport d’analyse forensique complet sur les téléphones & Extraction logique et physique (IMEI, contacts, messages, logs, métadonnées GPS). \\ \hline
Analyse des serveurs internes de la DGRE & Vérification de la traçabilité des ordres de mission et des connexions au réseau interne. \\ \hline
Corrélation temporelle automatisée (timeline numérique) & Reconstitution minute par minute des positions et communications des suspects. \\ \hline
Rapport de métadonnées sur la vidéosurveillance & Authentification (horodatage, intégrité des fichiers). \\ \hline
Analyse de l’utilisation du véhicule PRADO & Vérification du trajet effectué lors de l’enlèvement. \\ \hline
Extraction des transactions mobiles (Orange Money, Mobile Money) & Suivi des flux financiers entre acteurs avant/après le crime. \\ \hline
Analyse des messageries cryptées (Signal, Telegram, WhatsApp Business) & Détection de coordination clandestine via messagerie chiffrée. \\ \hline
Rapport d’expertise sur l’intégrité des données électroniques transmises & Garantir la validité juridique des preuves numériques. \\ \hline
\end{tabular}

\section{Synthèse des apports de l’investigation numérique}
\begin{itemize}[leftmargin=2em]
\item Les preuves électroniques ont permis de lier directement les auteurs aux lieux et aux moments précis des crimes.
\item L’enquête a reposé sur une corrélation entre données techniques et témoignages humains.
\item Les éléments numériques ont démonté les alibis, confirmé la préméditation et identifié la chaîne de commandement.
\item L’expert judiciaire a donc joué un rôle clé dans la matérialisation des infractions et la reconstitution complète du scénario criminel.
\end{itemize}

\section{Conclusion}
L’ordonnance de renvoi du 29 février 2024 repose sur une architecture numérique d’investigation complexe : exploitation des téléphones, géolocalisation, vidéosurveillance, données bancaires et métadonnées. 

Le travail d’un expert judiciaire en investigation numérique a été indispensable pour :
\begin{itemize}
\item consolider les preuves techniques,
\item établir la traçabilité des communications,
\item authentifier les éléments électroniques,
\item et permettre au magistrat d’instruire l’affaire jusqu’au renvoi devant la juridiction militaire.
\end{itemize}

\end{document}
